\begin{table}[H]
\centering\footnotesize
\setlength{\tabcolsep}{4pt}
\caption{Resultado orçamentário com e sem a dedução para o FUNDEB (R\$ milhões correntes)}
\label{tab:resultado}
\begin{adjustbox}{max width=\textwidth}
\begin{tabular}{lrrrrrr}
\toprule
 & \multicolumn{3}{c}{Barreiras} & \multicolumn{3}{c}{Camaçari} \\
\cmidrule(lr){2-4}\cmidrule(lr){5-7}
Ano & Bruto & Dedução & Líquido & Bruto & Dedução & Líquido \\
\midrule
2014 & 17,2 & 27,9 & -10,7$^{\dagger}$ & 81,9 & 88,9 & -7,0$^{\dagger}$ \\
2015 & 19,9 & 30,6 & -10,7$^{\dagger}$ & 80,2 & 98,4 & -18,2$^{\dagger}$ \\
2016 & 43,5 & 32,4 & 11,2 & 189,7 & 104,2 & 85,5 \\
2017 & 212,4 & 32,9 & 179,5 & 218,4 & 103,6 & 114,9 \\
2018 & -40,6 & 35,8 & -76,4 & 151,8 & 114,7 & 37,1 \\
2019 & -23,6 & 38,7 & -62,2 & 172,7 & 126,9 & 45,8 \\
2020 & 32,1 & 39,9 & -7,7$^{\dagger}$ & 227,1 & 124,4 & 102,6 \\
2021 & 74,1 & 52,3 & 21,8 & 184,4 & 145,2 & 39,2 \\
2022 & 20,7 & 63,7 & -43,1$^{\dagger}$ & 196,4 & 155,5 & 40,9 \\
2023 & 56,0 & 69,8 & -13,8$^{\dagger}$ & 288,4 & 157,8 & 130,6 \\
2024 & 88,0 & 81,8 & 6,2 & 176,6 & 157,6 & 19,0 \\
2025 & 47,8 & 89,3 & -41,6$^{\dagger}$ & 244,4 & 167,8 & 76,6 \\
\bottomrule
\end{tabular}
\end{adjustbox}
\par\vspace{3pt}\parbox{\textwidth}{\scriptsize Nota: Bruto = receita total bruta menos despesa empenhada, como no trabalho anterior; Líquido = receita total líquida da dedução para a formação do FUNDEB menos despesa empenhada. $\dagger$ = superávit no cálculo bruto que vira déficit no líquido. Entre 2014 e 2023, o resultado líquido é negativo em 7 anos em Barreiras (2014, 2015, 2018, 2019, 2020, 2022, 2023) e em 2 em Camaçari (2014, 2015). O resultado da execução de um exercício não mede, isoladamente, a situação fiscal: não considera restos a pagar, disponibilidade de caixa nem superávits financeiros de exercícios anteriores. Fonte: elaboração própria com dados do Siconfi (DCA).}
\end{table}
